% Part: normal-modal-logic
% Chapter: completeness
% Section: introduction

\documentclass[../../../include/open-logic-section]{subfiles}

\begin{document}

\olfileid{nml}{com}{int}

\olsection{Pambuka}

Yen $\Sigma$ sistem modal, teorema kasahihan netepake yen
$\Sigma \Proves !A$, mula $!A$ valid ing sembarang kelas
$\mClass{C}$ model kang kabeh instansi saka kabeh !!{formula}s
ing~$\Sigma$ valid. Mligine, yen $\Log{K}
\Proves !A$, mula $!A$ bener ing kabeh model; yen $\Log{KT} \Proves !A$,
mula $!A$ bener ing kabeh model refleksif; yen $\Log{KD} \Proves !A$,
mula $!A$ bener ing kabeh model serial, lan sateruse.

Kelengkapan iku kosok baline kasahihan: \Log{K} jangkep
ateges yen !!a{formula}~$!A$ valid, mula $\Proves !A$, upamane.
Mbuktekake kelengkapan luwih angel tinimbang mbuktekake kasahihan.
Mula, migunani yen luwih dhisik nimbang kontraposisine: \Log{K}
jangkep yen lan mung yen saben $\Proves/ !A$ ana model kang mbantah,
yaiku model~$\mModel{M}$ saengga $\mSat/{M}{!A}$. Kanthi cara
ekuivalen (kanthi negasi $!A$), kita bisa mbuktekake yen saben
$\Proves/ \lnot !A$ ana model saka~$!A$. Nalika mbangun model kaya
mangkono, kita bisa nggunakake informasi kang ana ing $!A$. Nalika
nemokake model kanggo !!{formula}s tartamtu, kita kerep nindakake bab
kang padha: upamane, yen arep nemokake model kang mbantah
$p \lif \Box q$, kita ngerti yen model iku kudu ngandhut jagad kang
$p$ bener lan $\Box q$ salah. Jagad kang $\Box q$ salah mbutuhake
anane jagad kang bisa diakses saka kono lan $q$ salah. Iku kabeh kang
perlu dingerteni: !!{propositional variable}s endi kang bener ing saben
jagad, lan jagad endi kang bisa diakses saka jagad liyane.

Nanging, nalika mbuktekake kelengkapan, kita ora duwe
!!{formula}~$!A$ tartamtu kang dibanguni model. Kita arep netepake
anane model kanggo saben $!A$ saengga
$\Proves/[\Sigma] \lnot !A$. Iki sarat minimal, amarga
\emph{yen} $\Proves[\Sigma] \lnot !A$, miturut kasahihan, ora ana model
kanggo~$!A$ (kang nyukupi $\Sigma$). Saiki elinga yen $\Proves/[\Sigma]
\lnot !A$ yen lan mung yen $!A$ iku $\Sigma$-konsisten. (Elinga yen
$\Proves/[\Sigma] \lnot !A$ lan $!A \Proves/[\Sigma] \lfalse$ iku
ekuivalen.) Dadi, tugas kita yaiku mbangun model kanggo saben
!!{formula} kang $\Sigma$-konsisten.

Carane yaiku nemokake himpunan !!{formula}s kang
$\Sigma$-konsisten lan ngandhut~$!A$, nanging uga ngemot formula
liyane kang nerangake wujude jagad kang ndadekake $!A$ bener. Himpunan
kaya mangkono iku \emph{jangkep} $\Sigma$-konsisten. Ora cukup yen
model kang ndadekake~$!A$ bener mung nduweni siji jagad; model iku
kudu nduweni luwih saka siji jagad lan relasi aksesibilitas. Himpunan
$\Sigma$-konsisten jangkep kang ngandhut~$!A$ uga ngandhut
!!{formula}s liyane awujud $\Box !B$ lan~$\Diamond !C$. Ing kabeh
jagad kang bisa diakses, $!B$ kudu bener; ing paling ora siji jagad,
$!C$ kudu bener. Kanggo nggayuh iki, kita njupuk \emph{kabeh}
himpunan $\Sigma$-konsisten jangkep kang bisa ana minangka dhasar
himpunan jagad. Bab kang angel yaiku nemtokake kapan sawijining
himpunan $\Sigma$-konsisten jangkep dianggep bisa diakses saka
himpunan liyane ing model kita.

Kita bakal nuduhake yen ing model kang ditepakake mangkono,
$!A$ bener ing sawijining jagad---kang uga himpunan $\Sigma$-konsisten
jangkep---yen lan mung yen $!A$ iku !!a{element} saka himpunan
kasebut. Yen $!A$ $\Sigma$-konsisten, formula iku bakal dadi
!!a{element} saka paling ora siji himpunan $\Sigma$-konsisten jangkep
(kasunyatan kang bakal kita buktekake), mula bakal ana jagad kang
$!A$~bener. Dadi, kita bakal nduweni siji model kang saben
$\Sigma$-konsisten !!{formula}~$!A$ bener ing sawijining jagad.
Model tunggal iki yaiku model \emph{kanonik} kanggo~$\Sigma$.

\end{document}
